% ─────────────────────────────────────────────────────────────────────────────
% Experimental Design
% Now a \subsection of the Methods section (fits the fixed thesis structure).
%
% Hyperparameter provenance: benchmark-specific values deferred to the
% appendix table.
% ─────────────────────────────────────────────────────────────────────────────

\subsection{Experimental Design}
\label{sec:experimental_design}

The four sub-research questions demand qualitatively different experimental
designs. Sub-RQ1 and Sub-RQ3 require head-to-head benchmark comparisons ---
projection methods against an unconstrained baseline, and \texttt{iFOPNG}
against \texttt{FOPNG} respectively. Sub-RQ2 requires a controlled ablation
over normalization conditions. Sub-RQ4 requires a long-horizon comparison of
two Fisher accumulation operators on the same inertial combined-metric preconditioner.

% ─────────────────────────────────────────────────────────────────────────────
\subsubsection{Hyperparameter Provenance}
\label{sec:hparam_provenance}

\paragraph{Standalone target-network experiments.}
All hyperparameter values for experiments without a hypernetwork ---
learning rates, regularization coefficients, gradient budgets, Fisher
sample counts, batch sizes, and epoch counts --- are taken without
modification from \citet{garg_fisher-orthogonal_2026} Table~1.
\texttt{iFOPNG} uses the same values as \texttt{FOPNG}, since the inertial
combined Fisher introduces no hyperparameters beyond those FOPNG already
defines.
\texttt{ONG} \citep{yadav_ong_2025} is treated as a composite
method: it applies \texttt{OGD}'s Euclidean projection operator to the
natural gradient direction and introduces no hyperparameters beyond those
already specified for \texttt{OGD} and \texttt{FNG}; its values therefore
follow \texttt{OGD}'s Table~1 entries.

\paragraph{Deviation for first-task initialization.}
We deviated from the original protocol in one specific instance: first-task 
initialization for MNIST-based benchmarks. While \citet{garg_fisher-orthogonal_2026} 
state they used the method's own learning rate for first-task SGD training 
(e.g., $\eta = 10^{-5}$ for FOPNG), our empirical replication of this protocol 
yielded an initial accuracy of only $\approx 21\%$. As demonstrated by the 
learning rate convergence sweep detailed in Appendix~D (Figure~\ref{fig:sgd_threshold}), 
this published rate falls well below the empirical convergence threshold, 
effectively failing to establish a functional baseline. To ensure a clean initial 
projection subspace, we standardized the first-task SGD learning rate to $10^{-3}$ 
for MNIST-based benchmarks and $10^{-2}$ for CIFAR-based benchmarks. 
Beyond this necessary empirical correction, no tuning was performed on the published values.

\paragraph{Hypernetwork experiments.}
Settings not covered by \citet{garg_fisher-orthogonal_2026} follow the canonical 
hypernetwork framework of \citet{oswald_continual_2020}; exact values are detailed 
in Appendix~\ref{app:hyperparameters}. Because computing a full Fisher estimate 
requires a complete $K$-chunk generation pass per training sample, estimation is 
capped at $1{,}024$ samples to maintain computational tractability --- a standard 
approximation consistent with prior CIFAR-10 evaluations \citep{kirkpatrick_overcoming_2017}. 
The gradient budget remains $K = 80$ per task, with memory compression handled via 
QR-and-SVD (Section~\ref{sec:projection_methods}). First-task training uses 
\texttt{AdamW} to seed the memory matrix $G$ and Fisher estimate $\hat{F}$ from a 
clean geometric basis.

\paragraph{Uniform learning rate control.}
To isolate the projection constraint as the sole independent variable in the 
hypernetwork experiments, all methods use a uniform learning rate of $10^{-3}$ 
rather than method-specific rates. This prevents convergence-speed artifacts 
from conflating the structural effects of the projection geometry.

\paragraph{Early stopping.}
Per-task training duration is governed by a single early-stopping rule
shared verbatim by every method, baseline and projection alike.
Training on a task terminates at the first of three conditions:
the average training loss falls below $0.1$;
ten consecutive epochs elapse without improving the best training loss
observed on that task;
or the benchmark's epoch cap (Appendix~\ref{app:hyperparameters}) is
reached.
A coupled scheduler reduces the learning rate after every three
consecutive non-improving epochs, stepping through half-magnitude
decrements ($10^{-2} \to 5 \times 10^{-3} \to 10^{-3} \to \dots$).
Upon termination, the parameters of the best-loss epoch --- not the final
epoch --- are restored and carried into the next task.
The criterion consults only the training loss, so the stopping rule never
observes held-out data; per-task duration is therefore adaptive while
introducing no method-specific degree of freedom.

\paragraph{Random seeds.}
All experiments draw from a fixed seed pool.
Configurations evaluated on five seeds use $\{42, 111, 811, 1234, 2137\}$;
compute-limited configurations use the three-seed subset
$\{42, 811, 1234\}$.
The seed count for every reported comparison is stated alongside the
corresponding result.
One deviation applies: on the Split-MNIST hypernetwork benchmark, seed $2137$ produced degenerate runs for every
projection-based method while leaving the unconstrained baselines
unaffected; it was therefore excluded and replaced by seed $314$,
yielding the five-seed set $\{42, 111, 314, 811, 1234\}$ used uniformly
across all methods on this benchmark.

\paragraph{Fisher accumulation default.}
The EMA decay parameter $\alpha$ appears exclusively in the Sub-RQ4
operator comparison ($\alpha = 0.5$); every other experiment accumulates
$\hat{F}_{\mathrm{old}}$ with the element-wise \textsc{max} operator,
which introduces no additional hyperparameter.

% 
% ─────────────────────────────────────────────────────────────────────────────
\subsubsection{Sub-RQ1 --- Constructive Synergy or Architectural Conflict}
\label{sec:design_rq1}

\paragraph{Experimental design.}
Sub-RQ1 investigates whether layering a hard geometric projection on top of a 
hypernetwork's soft functional regularizer yields a net benefit, or if the 
loss of adaptive momentum makes a standard \texttt{Adam}-regularized hypernetwork 
superior. To separate architectural conflicts from benchmark-specific artifacts, 
we evaluate across three datasets at varying levels of compression, comparing 
an unconstrained standalone network (Panel a) against a hypernetwork (Panel b).

\paragraph{Benchmark selection.}
We utilize a three-benchmark progression. \textbf{Benchmark 1 (Split-MNIST)} 
utilizes a deliberately ``suffocated'' hypernetwork (generator hidden dimension 
$d_h = 8$) to expose differential effects without hitting the near-perfect accuracy 
ceilings observed at higher capacities (see Appendix~\ref{app:hn_ablations}, Figure~\ref{fig:dh_sweep}). To ensure convergence under these constraints, each task is trained for 
$15$ epochs with a batch size of $64$. \textbf{Benchmark 2 (Split-CIFAR10)} serves 
as the primary test, offering genuine inter-task interference and utilizing a 
standard-capacity hypernetwork trained for 50 epochs. \textbf{Benchmark 3 (Split-CIFAR100)} 
extends this sequence to 10 tasks to test scalability. Comparing performance 
within and across these panels reveals whether projection methods universally 
dominate, conditionally succeed based on compression, or systematically fail 
due to momentum stripping.

% ─────────────────────────────────────────────────────────────────────────────
\subsubsection{Sub-RQ2 --- Projection-Scoped Normalization} % I SHOULD REWRITE THIS SO I EXPLAIN WHY A CONDITION IS A GOOD MEASURMENT TO EXPLAIN A inversion failures.
\label{sec:design_rq2}

\paragraph{What we want to differentiate.}
Sub-RQ2 asks whether the chunk-induced gradient accumulation pathology is
real, characterisable, and resolvable.
A single accuracy number cannot answer this: we need direct numerical
evidence that the pathology occurs without normalization, that normalization
resolves it, and that the resolution does not compromise projection quality.

\paragraph{Experimental design.}
\texttt{iFOPNG} is evaluated under three normalization conditions on the
Split-CIFAR10 hypernetwork (Benchmark~2), with all other hyperparameters
held constant.
Split-CIFAR10 is chosen over the CIFAR-100 and MNIST hypernetwork
configurations specifically because its chunk size of 6\,000 yields
$C = 198$ chunks per target network --- large enough to produce a
measurable pathology, yet small enough that the condition number of $A$
can be computed and logged at every projection step without becoming
numerically intractable.
Larger chunk counts (e.g.\ the CIFAR-100 configuration) would make the
condition number measurement itself unreliable; smaller counts (e.g.\
the MNIST suffocated regime) would not stress the pathology sufficiently.

\begin{description}
    \item[Condition 1 --- No normalization (negative control).]
    Incoming gradient batches are appended to $G$ without orthonormalization;
    the Fisher diagonal is used at its raw scale.
    With 198 chunks accumulated per forward pass, the gradient columns of % MAY BE REDUNDANT 
    $G$ and the entries of $\hat{F}$ are expected to reach extreme
    magnitudes, inflating the condition number of the projection kernel and
    driving the matrix inversion to \texttt{NaN}.
    This condition is the negative control that demonstrates the pathology
    concretely.

    \item[Condition 2 --- Gradient orthonormalization only.]
    Each incoming gradient batch is orthonormalized via QR decomposition
    before insertion into $G$; the Fisher diagonal is used at its raw scale.
    This isolates whether gradient-side conditioning alone is sufficient,
    or whether Fisher-side scaling is also required.

    \item[Condition 3 --- Full normalization (proposed).]
    Gradient batches are QR-orthonormalized at insertion, and the Fisher
    terms are rescaled by the maximum entry of the ambient metric $F^*$
    within each projection computation.
    This is the proposed scheme; all other hypernetwork experiments use it.
\end{description}

In addition to accuracy and BWT, the condition number $\kappa(A)$ of
the projection kernel is logged before each inversion.
A condition number that diverges in Condition~1 but remains bounded in
Condition~3 constitutes direct numerical evidence for the pathology and
its resolution.
All conditions use the Benchmark~2 hypernetwork settings, so any performance difference is attributable
to normalization alone.

% ─────────────────────────────────────────────────────────────────────────────
\subsubsection{Sub-RQ3 --- Does the Combined Fisher Metric Yield Consistent Retention Improvements}
\label{sec:design_rq3}

\paragraph{Benchmark selection and rationale.}
Sub-RQ3 compares \texttt{iFOPNG} against \texttt{FOPNG} across the four
standalone benchmarks evaluated by \citet{garg_fisher-orthogonal_2026}:
Permuted-MNIST (5 tasks), Split-MNIST single-head, Split-CIFAR10
multi-head, and Split-CIFAR100 multi-head.
All hyperparameters are taken directly from \citeauthor{garg_fisher-orthogonal_2026}'s
Table~1, with the sole exception of the first-task learning rate for
MNIST-based benchmarks, where the published value produces a degenerate
initialisation; the necessary correction is established empirically in
Appendix~D (Figure~\ref{fig:sgd_threshold}).
This replication-faithful design ensures that any observed difference
between \texttt{iFOPNG} and \texttt{FOPNG} is attributable to the combined
metric $F_c = \hat{F}_\mathrm{new} + \hat{F}_\mathrm{old}$ rather than
to differences in training conditions.

Split-MNIST multi-head and the two hypernetwork configurations
(Split-MNIST HN and Split-CIFAR10 HN) were already conducted for
Sub-RQ1 and are included here for completeness; they are not part of the
Garg et al.\ replication and carry no additional inferential weight for
Sub-RQ3.
% ─────────────────────────────────────────────────────────────────────────────

\subsubsection{Sub-RQ4 --- Fisher Accumulation Operator}
\label{sec:design_rq4}

\paragraph{What we want to differentiate.}
The non-decaying maximum $F_{\mathrm{old}} \leftarrow
\max(F_{\mathrm{old}}, F_{\mathrm{new}})$ guarantees that any parameter
ever deemed important to a past task retains that importance permanently.
This is protective for early tasks but monotonic by construction: the
inertial metric boundary can only tighten, never relax, so in a
sufficiently long sequence the accumulated metric may immobilise the
parameters needed to learn later tasks.

\paragraph{Experimental design.}
\texttt{iFOPNG} is run under two accumulation operators, all other settings
held constant:
\begin{description}
    \item[Maximum (proposed).] $F_{\mathrm{old}} \leftarrow
    \max(F_{\mathrm{old}}, F_{\mathrm{new}})$ --- non-decaying.
    \item[Exponential moving average.] $F_{\mathrm{old}} \leftarrow
    (1-\alpha)\,F_{\mathrm{old}} + \alpha\,F_{\mathrm{new}}$ ---
    decaying, allowing historical importance to fade.
\end{description}
The comparison is run on the longest available task sequences ---
Permuted-MNIST --- since the plasticity cost of
monotonic accumulation, if present, only becomes observable as the
sequence length grows.
Per-task accuracy is tracked across the full sequence rather than reported
only as a final aggregate: the diagnostic signature of plasticity paralysis
is a progressive decline in \emph{newly learned} task accuracy
($a_{t,t}$ falling as $t$ increases), distinct from forgetting, which
appears as decline in \emph{previously learned} task accuracy.
If the maximum operator shows superior early-task retention but degrading
$a_{t,t}$ in late tasks, while the moving average shows the opposite trade,
the two failure modes are cleanly separated and Sub-RQ4 is answered in
full.